% ==========================================================
%  PREÁMBULO — Ecuaciones Diferenciales, Tema 2
%  Motor: pdfLaTeX · Fuente: Palatino (newpx), como el artículo de referencia
% ==========================================================
\usepackage[utf8]{inputenc}
\usepackage[T1]{fontenc}
\usepackage[spanish,es-tabla,es-nodecimaldot]{babel}
\usepackage[p,osf]{newpxtext}
\usepackage{amsmath, amsthm, mathtools}
\usepackage[bigdelims,vvarbb]{newpxmath}
\usepackage[final,protrusion=true,expansion=true]{microtype}
\usepackage{bm}
\usepackage{siunitx}
\usepackage{graphicx}
\usepackage{enumitem}
\usepackage{booktabs}
\usepackage{array}
\usepackage{float}
\usepackage{fancyhdr}
\usepackage{needspace}
\usepackage{titlesec}
\usepackage{xcolor}
\usepackage{caption}
\usepackage[a4paper,left=3cm,right=3cm,top=3cm,bottom=3cm]{geometry}

\sisetup{locale=FR, per-mode=symbol, inter-unit-product=\ensuremath{{}\cdot{}}}
\sisetup{output-decimal-marker={,}}

% --- Paleta sobria (azules y grises) ---
\definecolor{azulOscuro}{HTML}{1F3A5F}
\definecolor{azulMedio}{HTML}{3B6EA5}
\definecolor{azulClaro}{HTML}{8FB3D9}
\definecolor{azulFondo}{HTML}{F3F7FB}
\definecolor{gris}{HTML}{6B7280}
\definecolor{grisClaro}{HTML}{D4D8DE}
\definecolor{acento}{HTML}{B5651D}

% --- Enlaces e índice ---
\usepackage[colorlinks=true,linkcolor=azulOscuro,citecolor=azulOscuro,urlcolor=azulMedio]{hyperref}
\usepackage[nameinlink,noabbrev,spanish]{cleveref}
\usepackage{tocloft}
\renewcommand{\cfttoctitlefont}{\hfill\large\scshape\color{azulOscuro}}
\renewcommand{\cftaftertoctitle}{\hfill}
\renewcommand{\cftsecleader}{\hfill}
\renewcommand{\cftsubsecleader}{\hfill}
\renewcommand{\cftsecfont}{\color{azulMedio}}
\renewcommand{\cftsubsecfont}{\color{azulMedio}}
\renewcommand{\cftsecpagefont}{\color{azulMedio}}
\renewcommand{\cftsubsecpagefont}{\color{azulMedio}}
\setcounter{tocdepth}{2}

% --- Tipografía y página ---
\linespread{1.08}
\setlength{\parindent}{1.2em}
\setlength{\parskip}{0pt}
\setlist{itemsep=2pt, topsep=4pt}
\pagestyle{fancy}
\fancyhf{}
\renewcommand{\headrulewidth}{0pt}
\setlength{\headheight}{14pt}
\fancyhead[C]{\footnotesize\scshape\color{gris} Ecuaciones diferenciales · Tema 2}
\fancyhead[R]{\footnotesize\color{gris}\thepage}
\fancypagestyle{plain}{\fancyhf{}\fancyfoot[C]{\footnotesize\color{gris}\thepage}\renewcommand{\headrulewidth}{0pt}}

% Estilo AMS: sección centrada en versalitas; subsección en negrita
\titleformat{\section}[block]{\centering\scshape\color{azulOscuro}}{\thesection.}{0.5em}{}
\titleformat{\subsection}[block]{\bfseries\color{azulOscuro}}{\thesubsection.}{0.5em}{}
\titleformat{name=\subsection,numberless}[block]{\itshape\color{azulOscuro}}{}{0pt}{}
\titlespacing*{\section}{0pt}{3ex plus 1ex}{1.6ex}
\titlespacing*{\subsection}{0pt}{2.2ex plus .8ex}{1ex}

% --- Notación ---
\newcommand{\dd}{\mathrm{d}}
\newcommand{\ee}{\mathrm{e}}
\newcommand{\dv}[2]{\dfrac{\dd #1}{\dd #2}}
\newcommand{\R}{\mathbb{R}}
\newcommand{\abs}[1]{\lvert #1 \rvert}
\newcommand{\eqnum}[1]{\tag{#1}}
\mathtoolsset{showonlyrefs=false}
\numberwithin{equation}{section}

% --- Gráficas ---
\usepackage{tikz}
\usetikzlibrary{arrows.meta, calc, positioning, decorations.markings}
\usepackage[siunitx]{circuitikz}
\usepackage{pgfplots}
\pgfplotsset{compat=1.17}
\pgfplotsset{
  informe/.style={
    width=0.9\linewidth, height=6.4cm,
    axis lines=left, axis line style={gris},
    tick style={gris}, tick label style={font=\footnotesize},
    label style={font=\small},
    grid=major, grid style={grisClaro, very thin},
    enlargelimits=false, clip=true,
    legend style={font=\footnotesize, draw=none, fill=none, at={(0.97,0.5)}, anchor=east, legend cell align=left},
  },
  serie1/.style={color=azulOscuro, line width=1.1pt},
  serie2/.style={color=azulMedio, line width=1.1pt, dashed},
  serie3/.style={color=gris, line width=1.1pt, dotted},
  serie4/.style={color=acento, line width=1.1pt, dash dot},
  asintota/.style={color=gris, line width=0.6pt, dashed},
}
\captionsetup{font=small, labelfont={bf,color=azulOscuro}, labelsep=period}

% ==========================================================
%  ENTORNOS
% ==========================================================
\usepackage[most]{tcolorbox}

\newcommand{\etiquetaEntorno}[1]{\textsc{#1}}

% Teoremas, proposiciones, lemas, corolarios: cuerpo en cursiva, encabezado azul
\newtheoremstyle{enunciado}{9pt}{9pt}{\itshape}{}{\bfseries\color{azulOscuro}}{.}{0.6em}%
  {\thmname{#1}\ \thmnumber{#2}\thmnote{\ (#3)}}
% Ejemplos y observaciones: cuerpo recto, encabezado en cursiva azul
\newtheoremstyle{nota}{8pt}{8pt}{\normalfont}{}{\itshape\bfseries\color{azulMedio}}{.}{0.6em}%
  {\thmname{#1}\ \thmnumber{#2}\thmnote{\ (#3)}}

\theoremstyle{enunciado}
\newtheorem{theorem}{Teorema}[section]
\newtheorem{proposition}[theorem]{Proposición}
\newtheorem{lemma}[theorem]{Lema}
\newtheorem{corollary}[theorem]{Corolario}
\theoremstyle{nota}
\newtheorem{example}{Ejemplo}[section]
\newtheorem{remark}[example]{Observación}

\renewcommand{\qedsymbol}{\textcolor{azulMedio}{$\blacksquare$}}
\renewcommand{\proofname}{\itshape Demostración}
\newcommand{\finejemplo}{\hfill\textcolor{azulMedio}{$\diamond$}}

% Definiciones: único entorno enmarcado (fino, fondo tenue)
\newtcbtheorem[number within=section]{definition}{Definición}{
  enhanced, breakable, sharp corners,
  colback=azulFondo, colframe=azulClaro, boxrule=0.4pt,
  left=9pt, right=9pt, top=5pt, bottom=5pt,
  before skip=10pt, after skip=10pt,
  fonttitle=\bfseries\color{azulOscuro}, coltitle=azulOscuro,
  attach boxed title to top left={yshift*=-\tcboxedtitleheight/2, xshift=8pt},
  boxed title style={colback=white, colframe=white, sharp corners, boxrule=0pt, left=3pt, right=3pt, top=1pt, bottom=1pt},
  title={}}{def}
\newenvironment{definicion}[2][]{\begin{definition}{#2}{#1}}{\end{definition}}

% Modelo: filete superior e inferior
\newtcolorbox{modelo}[1][]{
  enhanced, breakable, blanker,
  borderline north={0.6pt}{0pt}{azulOscuro},
  borderline south={0.6pt}{0pt}{azulOscuro},
  top=6pt, bottom=6pt, left=2pt, right=2pt,
  before skip=10pt, after skip=10pt,
  before upper={\noindent{\scshape\bfseries\color{azulOscuro} Modelo\if\relax\detokenize{#1}\relax\else\ --- #1\fi.}\ }}

% Desarrollo teórico: línea lateral fina gris
\newtcolorbox{desarrollo}[1][]{
  enhanced, breakable, blanker,
  borderline west={0.5pt}{0pt}{grisClaro},
  left=12pt, top=2pt, bottom=2pt,
  before skip=8pt, after skip=8pt,
  before upper={\noindent{\scshape\color{gris} Desarrollo\if\relax\detokenize{#1}\relax\else\ --- #1\fi.}\ }}

% Algoritmo de resolución: pasos numerados
\newcounter{algoritmo}[section]
\renewcommand{\thealgoritmo}{\thesection.\arabic{algoritmo}}
\newtcolorbox{algoritmo}[1][]{
  enhanced, breakable, blanker,
  borderline north={0.6pt}{0pt}{azulOscuro},
  borderline south={0.3pt}{0pt}{azulClaro},
  top=6pt, bottom=6pt, left=2pt, right=2pt,
  before skip=10pt, after skip=10pt,
  before upper={\refstepcounter{algoritmo}\noindent{\scshape\bfseries\color{azulOscuro} Algoritmo \thealgoritmo\if\relax\detokenize{#1}\relax\else\ (#1)\fi.}\par\vspace{2pt}},
  }
\newlist{pasos}{enumerate}{1}
\setlist[pasos]{label=\textbf{\textcolor{azulMedio}{Paso \arabic*.}}, leftmargin=3.6em, labelsep=0.5em, itemsep=2pt}

% Ejercicios
\newcounter{ejercicio}[section]
\renewcommand{\theejercicio}{\thesection.\arabic{ejercicio}}
\newtcolorbox{ejercicio}[1][]{
  enhanced, breakable, blanker,
  borderline west={1.2pt}{0pt}{azulMedio},
  left=12pt, top=3pt, bottom=3pt,
  before skip=10pt, after skip=10pt,
  before upper={\refstepcounter{ejercicio}\noindent{\bfseries\color{azulOscuro} Ejercicio \theejercicio\if\relax\detokenize{#1}\relax\else\ (#1)\fi.}\ }}

% Tabla de parámetros
\newcommand{\encabezadoTabla}[1]{\textbf{\textcolor{azulOscuro}{#1}}}

\crefname{figure}{figura}{figuras}
\Crefname{figure}{Figura}{Figuras}
\crefname{equation}{ecuación}{ecuaciones}
\Crefname{equation}{Ecuación}{Ecuaciones}
\crefname{theorem}{teorema}{teoremas}
\crefname{proposition}{proposición}{proposiciones}
\crefname{lemma}{lema}{lemas}
\crefname{algoritmo}{algoritmo}{algoritmos}

% Resumen estilo AMS
\newenvironment{resumen}{\small\begin{quote}\noindent\textsc{Resumen.}\ }{\end{quote}}

% --- Utilidades adicionales ---
\usepackage{subcaption}
\captionsetup[subfigure]{font=footnotesize, labelfont={color=azulOscuro}, labelsep=period}
\newcommand{\metodo}[1]{\par\smallskip\noindent\textbf{\textcolor{azulMedio}{#1.}}\ }
\newenvironment{solucion}{\par\smallskip\noindent\textit{\textcolor{azulOscuro}{Solución.}}\ }{\hfill\textcolor{azulMedio}{$\diamond$}\par\smallskip}
\newcommand{\Tf}{t_{\mathrm{f}}}

% Problemas de seminario
\newcounter{problema}[section]
\renewcommand{\theproblema}{\thesection.\arabic{problema}}
\newtcolorbox{problema}[1][]{
  enhanced, breakable, blanker,
  borderline west={1.2pt}{0pt}{azulMedio},
  left=12pt, top=3pt, bottom=3pt,
  before skip=12pt, after skip=6pt,
  before upper={\refstepcounter{problema}\noindent{\bfseries\color{azulOscuro} Problema \theproblema\if\relax\detokenize{#1}\relax\else\ (#1)\fi.}\ }}
